%   ___                       _              
%  / __|___ _ _  ___ _ _ __ _| |_ ___ _ _ ___
% | (_ / -_) ' \/ -_) '_/ _` |  _/ _ \ '_(_-<
%  \___\___|_||_\___|_| \__,_|\__\___/_| /__/
%

\chapter{Generators}

\noindent The rest of this document concerns the relationships between Lie
groups and Lie algebras. The first of these relationships comes via 
\firstmention{generators}.

\begin{definition}[Generator]
    A generator $X\in\mathfrak{g}$ is a Lie-algebra element that, multiplied by
    a real $\theta$,\sidenote{$X$ is a vector in the Lie Algebra $\mathfrak{g}$
    considered as a vector space, and, by the rules of vector spaces,
    vectors can always be multiplied by scalars from the underlying field}
produces a an element of a Lie group G under the (abstract) exponential map:
    \begin{align*}
        \exp &: (\mathfrak{g} \times \mathbb{R}) \rightarrow G \\
        h(\theta) &= \exp(\theta{}X)
    \end{align*}
\end{definition}

\begin{remark} A generator $X\in\mathfrak{g}$ in the Lie algebra can be recovered 
    from a Lie-group element $h(\theta)$ by (abstract) differentiation: 
    \begin{equation*}
        X = \left.\frac{d h(\theta)}{d \theta}\right|_{\theta=0}
    \end{equation*}
\end{remark}

\section{A representation is a concretization}

The exponential map above is abstract, as is the differentiation operation to
recover the Lie-algebra element. $\operatorname{Exp}$ must be further defined
for any concretization. For the sake of computation, physicists work almost
exclusively with matrix concretizations. In such concretizations, $\exp$ is a
matrix exponential, defined by Taylor series. From this point on, assume all
concretizations are matrix concretizations. 

Loosely, a \firstmention{representation} is a matrix that stands for a
Lie-algebra or Lie-group element. We formalize \textit{representation} later,
but for the purpose of illustrating generators as matrices, we must choose
representations.

\section{The same Lie algebra can generate many Lie Groups}

We will present some general results and then illustrate them with two
particular examples of interest in physics. Those two examples are
representations of one Lie algebra \su{2} that generate two different Lie
groups. In the first representation, the \firstmention{adjoint
representation},\sidenote{symbolised by \adjrep}\sidenote{no relation to the
adjoint operation on matrices, which is complex-conjugate transpose.} elements
of the Lie algebra are represented by $3\times{}3$ real matrices. In the second
representation, the \firstmention{fundamental
representation},\sidenote{symbolized by \fndrep} elements are represented by
$2\times{}2$ complex matrices. We will show how to map from one to the other
and back, that is, the \firstmention{isomorphism} between these two
representations. The meaning of that isomorphism is that the two
representations actually represent the same Lie algebra. However, the two
generated Lie groups are \textit{not} equivalent, and interesting physics
follow from that fact.

For the time being, take the words \textit{adjoint} and \textit{fundamental} at
face value. We will elaborate them extensively as we go along.

\section{Some general results}

\begin{proposition}
    \label{prop:traceless-antisymmetric}
    Any real matrix generators $J\in\mathfrak{g}$ of a \firstmention{special}
    \firstmention{orthogonal} group \SO{N} are traceless and antisymmetric.
\end{proposition}

\begin{remark}
    \SO{N} is a group of $N\times{}N$ invertible\sidenote{All elements of any
    group are invertible, by definition of a group} matrices with two more
    properties:
    \begin{itemize}
        \item \textit{Special} means that the elements of \SO{N} have unit determinant. 
        \item \textit{Orthogonal} means that the transpose $\transpose{O}$ of any
            element equals its matrix inverse $O^{-1}$.
    \end{itemize}
\end{remark}

\begin{remark}
    A generator in the Lie algebra, if it is to geneate an $N\times{}N$ real matrix in the Lie group by the
    matrix exponential, must itself be represented by an $N\times{}N$ real matrix.
\end{remark}

\begin{proof}
    The proof proceeds from implications of the definition of a matrix exponential.
\leavevmode
\begin{itemize}[noitemsep] % Removes space between items
    \item \textit{Traceless}:
        \textit{Special} implies that the $N$\nobreakdash-dimensional
        generators are traceless:\sidenote{Equation numbers with prefixes, such
        as (Sch:3.60), are taken from particular references that are meant to
        be obvious from the prefix.}
        \begin{equation}
            \begin{split}
                \det{O}\mustbe{1}&\quad\implies\quad\det{e^{\theta{}J}}=e^{\theta\;\trace{J}}=1 \\
                e^{\theta\;\trace{J}}=1&\implies\theta\;\trace{J}=0\implies\trace{J}=0
            \end{split}
            \tag{Sch:3.60}
        \end{equation}
        \noindent by \firstmention{Jacobi's formula} linking $\det{\exp}$ and 
        $\exp{(\operatorname{Tr})}$.\cite{horn2012matrix}
    \item \textit{Antisymmetric}: 
        \textit{Orthogonal} means that elements $O$ of \SO{3} satisfy 
        $\transpose{O}O=O\transpose{O}\mustbe{1}$. Therefore:
        \begin{equation}
            \begin{split}
            {\transpose{O}}O=e^{\theta\transpose{J}}
                               e^{\theta{}J}\mustbe{}1
                               &\quad\implies\quad
                           e^{\theta\;(\transpose{J}+J)}=1\\
                       &\quad\implies\quad 
                       \transpose{J}=-J
            \end{split}
        \tag{Sch:3.59}
    \end{equation}
    \noindent by the Baker-Campbell-Hausdorff formula\cite{hall2015lie} for matrix
    exponentials, because $\transpose{J}$ and $J$ commute.
    The \textit{transpose} operation lifts into the exponent, as can be see by
    a Taylor expansion, because $\theta$ is a scalar.
\end{itemize}
\end{proof}

\section{The Adjoint Representation is a particular real 3x3 representation}

\SO{3} is our controlling example: it is the Lie group of rotations in 3D
Euclidean space, familiar to all physics students. 

Proposition~\ref{prop:traceless-antisymmetric} pertains to any special
orthogonal group \SO{N}. How do we get members of \SO{3} from generators in
\su{2}? First, we state\sidenote{without proof, for now} that \su{2} is a
3\nobreakdash-dimensional vector space, so we can choose a basis of three
linearly independent generators $\{X_1, X_2, X_3\}$ to span \su{2}. 

\SO{3} is also three dimensional, in the different sense that it's full of
$3\times{}3$ matrices. 

The following basis makes up the ajoint representation of \su{2} that generates
\SO{3}. This set is not a guess; we'll show later how to derive it from the
fundamental representation.\sidenote{The Lie bracket can ``generate'' the basis
    (in a different sense of the word \textit{generate}), from just two of the
$J$ because $[J_1, J_2]=J_3$, etc.}

{
\small
\addtolength{\arraycolsep}{-3pt}
\begin{equation}
    J_1 = 
    \begin{pmatrix}
        0 & 0 & 0 \\
        0 & 0 & -1 \\
        0 & 1 & 0
    \end{pmatrix}, \quad
    J_2 = 
    \begin{pmatrix}
        0 & 0 & 1 \\
        0 & 0 & 0 \\
        -1 & 0 & 0
    \end{pmatrix}, \quad
    J_3 = 
    \begin{pmatrix}
        0 & -1 & 0 \\
        1 & 0 & 0 \\
        0 & 0 & 0
    \end{pmatrix}
    \tag{Sch:3.61}
\end{equation}
}

The following Mathematica code checks that these generators in \su{2} satisfy
the Lie-algebra commutation relations $[J_i, J_j] = \epsilon_{ijk}
J_k$:\sidenote{$\epsilon_{ijk}$ is the totally antisymmetric Levi-Civita symbol
    that equals $1$ when $ijk$ is an even permutation of $123$, $-1$ when
${ijk}$ is an odd permutation, and $0$ otherwise. This expression does
\textit{not} use the Einstein summation convention; it's just ordinary
scalar-vector (as matrix) multiplication for any $i$, $j$, and $k$.}

\begin{Verbatim}[frame=single, label=Mathematica]
J1 = {{0,  0, 0}, {0, 0, -1}, { 0, 1, 0}};
J2 = {{0,  0, 1}, {0, 0,  0}, {-1, 0, 0}};
J3 = {{0, -1, 0}, {1, 0,  0}, { 0, 0, 0}};

J1.J2 - J2.J1 == J3 (* True *)
J2.J3 - J3.J2 == J1 (* True *)
J3.J1 - J1.J3 == J2 (* True *)
\end{Verbatim}

\noindent Mathematica can further check that the matrix exponentials of these
generators span \SO{3} as a matrix group. The following should convince you
if you are familiar with elementary rotation matrices:

\begin{Verbatim}[frame=single, label=Mathematica]
MatrixExp[\[Theta] J1]
{{1, 0, 0}, 
 {0, Cos[\[Theta]], -Sin[\[Theta]]}, 
 {0, Sin[\[Theta]],  Cos[\[Theta]]}}
\end{Verbatim}

\section{The Fundamental Representation is a distinguished complex 2x2 representation}

\begin{fact}
There is exactly one simply connected Lie group corresponding to each Lie algebra.
\end{fact}

\begin{proof}
    See Spivak.\cite{spivak1999}
\end{proof}

\noindent We don't need yet to understand ``simply connected'' to
appreciate that the simply connected Lie group is the important and
distinguished one. For our controlling example, this group is \SU{2}, 
the special \firstmention{unitary} group of $2\times{}2$ matrices.

\begin{remark}\leavevmode
    \begin{itemize}
        \item \textit{Special} means that the elements of \SU{2} have unit determinant, as before.
        \item \textit{Unitary} means that the complex-conjugate transpose $\adjoint{U}$ of any
            element $U\in\SU{2}$ equals its matrix inverse $U^{-1}$. \textit{Unitary} for complex 
            matrices is the analog of \textit{orthogonal} for real matrices.
    \end{itemize}
\end{remark}

\begin{proposition}
    Complex $2\times{}2$ generators $T\in\su{2}$ of the special unitary group
    \SU{2} are traceless and
    \firstmention{anti-Hermitian}.\sidenote{\textit{Anti-Hermitian} for complex
    matrices is the analog of antisymmetric for real matrices.}
\end{proposition}

\begin{proof}

    Proof of the traceless property goes just as above. 

    \textit{Unitary} means that the elements $U$ of \SU{2} satisfy
    $\adjoint{U}U=U\adjoint{U}\mustbe{1}$,\sidenote{
        As noted, this meaning of \textit{adjoint} is distinct from its meaning 
        in \textit{adjoint representation}, which pertains to the representation of \SO{3}.}

    \textit{Hermitian} means that a matrix is self-adjoint. \textit{Anti-Hermitian} means
    that a matrix equals the negative of its adjoint. Now we see that
        \begin{equation}
            {\adjoint{U}}{U}=e^{\theta\adjoint{T}}
                             e^{\theta{}T}\mustbe{}1
                       \quad\implies\quad
                       \adjoint{T}\mustbe{}-T
        \end{equation}

\noindent The adjoint operation also lifts into the exponent, like
\textit{transpose}, because $\theta$ is a \textit{real} scalar. 

\end{proof}

\noindent A basis for this fundamental
representation is $-i/2$ times the three
\firstmention{Pauli matrices}:\sidenote{The Pauli matrices are also not a 
    guess. They are the unique $2\times{}2$ traceless Hermitian matrices
    in the complex domain. The factor of $-i$ makes them anti-Hermitian,
    and the factor of $1/2$ is convenient.}

{
\small
\addtolength{\arraycolsep}{-3pt}
\begin{equation}
    \sigma_1 = \begin{pmatrix} 0 & 1 \\ 1 & 0 \end{pmatrix}, \quad
    \sigma_2 = \begin{pmatrix} 0 & -i \\ i & 0 \end{pmatrix}, \quad
    \sigma_3 = \begin{pmatrix} 1 & 0 \\ 0 & -1 \end{pmatrix}.
    \tag{Sch:3.80}
\end{equation}
}

\noindent The basis is $T_1 = -i\sigma_1/2$, $T_2 = -i\sigma_2/2$, $T_3 = -i
\sigma_3/2$ as can be readily verified in Mathematica, showing, again in
passing, that the factor of $1/2$ maintains the unit \firstmention{structure
constants} $\epsilon_{ijk}$. Without the $1/2$, we'd be chasing around factors of $2$. 

\begin{Verbatim}[frame=single, label=Mathematica]
T1 = -I/2 PauliMatrix[1];
T2 = -I/2 PauliMatrix[2];
T3 = -I/2 PauliMatrix[3];
T1.T2 - T2.T1 == T3 (* True *)
T2.T3 - T3.T2 == T1 (* True *)
T3.T1 - T1.T3 == T2 (* True *)
\end{Verbatim}
